\subsubsection{路径生成结果概述}

P1 的实验重点是验证二维仓库场景中的候选路径生成是否稳定、完整且可复现。本文分别使用 A*、VG、AVG 和 DAVG 生成关键节点之间的基础路径源，并输出路径成本表与轨迹采样表。不同算法在路径长度、转弯数量和风险规避倾向上存在差异，因此能够用于比较不同路径生成机制的几何特征和代价结构。

需要说明的是，mixed 多候选路径库的融合、路径源对调度结果的影响，以及不同路径源下 $C_{\max}$、$F_2$、$F_3$ 的比较，均属于任务级调度实验，本文在 P2 实验章节中统一展开。P1 在此处只说明基础路径源的生成质量、算法差异和输出文件完整性。
